\section{Synchronous Systems}
Trattiamo ora i \textit{Fully Synchronous Systems}.

Clock e sincronizzati e delay limitati (bounded).
I clock sono sincronizzati e hanno periodo $\delta$.\\
Corollario: tutti i messaggi partono contemporaneamente.

Restrizioni: bounded communication delays. Esiste ed è conosciuto da ogni nodo un upper bound $\Delta$ sul delay di trasmissione.
I tempi di trasmissione non sono comunque costanti, ma ho la garanzia che ogni messaggio arrivi entro $\Delta$.

Posso riadattare l'unità di tempo a delta: $u = \lceil \frac{\Delta}{\delta} \rceil$, che rappresenta il massimo numero di tick necessari
per inviare qualsiasi messaggio. Consideriamo quindi il sistema sincrono con delay unitario (messaggi inviati e processati in al massimo
una time unit) pari a $u$. 

Sarebbe irragionevole avere un UB sul delay con messaggi di dimensione arbitraria, potenzialmente infinitamente grande.
Dobbiamo limitare la dimensione dei messaggi: i messaggi suono upper bounded.
Abbiamo un parametro di sistema $c$ che indica che ogni messaggio avrà dimensione massimo $c$ bits.
Messaggi più grandi devono essere spezzati per rientrare nella dimensione massima. 

Queste assunzioni semplificano il calcolo della time complexity siccome gli orologi sono tutti sincronizzati.
Data una distanza e un tempo globale, posso calcolare la velocità di un messaggio.

\subsection{Speeding}
Algoritmo di leader election su \underline{ring sincroni}.

Nota: gli algoritmi asincroni possono funzionare anche su sistemi sincroni, ma non il contrario.\\
Idea: come in \textit{As far as it can} ma in cui velicizziamo i messagi con gli ID più piccoli in
modo che i sconfiggano gli ID più grandi il prima possibile.\\
Restrizioni: tutti i nodi partono insieme, ring sincrono, unidirezionale. Non serve conoscenza di N numero di nodi.

Introduciamo un ritardo dipendente dall'ID, per differenziare le velocità.
Quando un nodo riceve un messaggio contenente $i$, se $i < x$, aspetta $f(i)$ unità di tempo, in caso contrario, se $i > x$, lo scarta.
Diventa leader il nodo che vede ritornarsi il prorprio ID; appena accade lo notifica (senza delay) alla rete.

Una funzione che funziona :) è l'esponenziale: $2^i$; gli ID maggiori aspetteranno più tempo, mentre gli ID minori li sorpasseranno.
Voglio che solo il leader riceva il messaggio che contiene il suo ID; nessun altro deve riceverlo prima di lui.

Quanto ci mette min ad attraversare il ring? $2^\text{min} * n + n$ ovvero il tempo di attesa più il tempo di attraversamento degli $n$ link.

Questo arriva prima di tutti gli altri? vediamo se il secondo più piccolo, ovvero $min+1$, fa in tempo a chiudere un giro.
Nel tempo che impiega il primo per chiudere un giro, riesce a percorrere al massimo $(2^\text{min} * n + n) / 2^{\text{min}+1} \leq n/2 + n/2 = n$ entità,
quindi strettamente meno di $n$ link.

Il secondo più piccolo riesce a percorrere al meno di $(2^\text{min} * n + n) / 2^{\text{min}+2} \leq n/4 + n/4 = n/2$ link.s

Il numero totale di messaggi è $\text{Messages}(\text{Speeding}(R)) = n + n/1 + n/2 + \dots + n/(2^{n-2}) \in O(n)$.
Il numero di bit è: $\text{Bits}(\text{Speeding}(R)) \in O(n \log(\text{max}))$.
L'attesa massima totale è $O(n * 2^\text{min})$, ovvero il tempo necessario al minimo per riuscire a tornare al leader.

\subsection{Waiting}
Il silenzio è espressivo. Possiamo usare il silenzio per comunicare informazioni.

Ad esempio, possiamo comunicare qualsiasi numero tra i due nodi con solo 2 messaggi da 1 bit l'uno, a patto che i clock siano molto
precisi. Dato un numero da comunicare $N$, ad un istante di tempo $t$ mando il bit di inizio (eg. un messaggio con 0) e dopo $N$ tick di clock
mando un messaggio di stop (eg. un messaggio con 1); il ricevente calcola la differenza tra i clock: $(t+N+1) - (t+1)$, ottenendo l'informazione.

Leader election su ring sincroni. \\
Idea: il leader aspetta un certo tempo prima di decidere che è stato eletto. \\
Restrizioni: come prima + \underline{conoscenza del numero di nodi n}, necessaria per sapere i tempi di comunicazione. \\

Reminder: su ring orientati il numero massimo di hop per collegare qualsiasi coppia di nodi è $n-1$; tra due nodi esiste un solo percorso. \\
Assunzione: tutti le entità iniziano simultaneamente. \\

La funzione di attesa dipende da $i$ e $n$: $f(i, n)$. Gli ID più piccoli aspettano meno, mentre gli ID più grandi aspettano di più. \\
Il vincolo che dobbiamo rispettare per il funzionamento del protocollo è che la notifica del leader raggiunga tutte le entità mentre
stanno ancora aspettando. Quello più distante per il messaggio di notifica è il nodo prima di lui, che ha distanza $n-1$
(ricorda che il ring è orientato); il caso più rischioso è quello in cui tale nodo ha ID pari a $x+1$.

Formalmente, siamo nel caso peggiore se leader $\text{ID} = x$; Dato il nodo con secondo più piccolo $y = x+1$; data la distanza tra i due
$d(x, y) = n-1$.
Deve essere vero che: $f(x+1, n) - f(x, n) >= n-1$. Ovvero, l'attesa del secondo minimo deve essere almeno il tempo necessario alla notifica
del leader per raggiungerlo, nel caso in cui si trovi nella posizione peggiore, ovvero appena ``dietro'' al leader.
La funzione $f(i, x) = i*x$ è la più piccola che soddisfa questo vincolo.

Questo funziona con \underline{inizio simultaneo}, tigliendo questo vincolo abbiamo bisogno di una nuova strategia basata su \textbf{wakeup}.
Ogni nodo sveglia il vicino e inizia ad aspettare. Vogliamo che il futuro leader finisca il waiting prima di tutti gli altri,
indipendentemente dal tempo di sveglia. Dobbiamo riprogettare la funzione di attesa.
Dato $t(j)$ tempo di sveglia di $j$: $f(i, n) \text{ s.t. } t(x) + f(x, n) + d(x, y) < t(y) + f(y, n)$.

Sicuramente:
\begin{itemize}
  \item $t(x) - t(y) < n$. Il tempo di svelia tra ogni coppia di nodi è al masismo $n-1$
  \item $d(x, y) < n$. La distanza tra due nodi è al max $n-1$.
  \item => $t(x) - t(y) + d(x, y) < 2n$
  \item => $t(x) - t(y) + d(x, y) + f(x, n) < 2n + f(x, n)$
  \item => $t(x) + d(x, y) + f(x, n) < 2n + f(x, n) + t(y)$
\end{itemize}

Ci chiediamo se: $t(x) + d(x, y) + f(x, n) < [2n + f(x, n)] + t(y) < t(y) + f(y, n)$

Cerchiamo una funzione tale che $2n + f(x, n) < f(y, n) \Rightarrow 2n < f(y, n) - f(x, n)$. Che nel caso peggiore ($y=x+1$) diventa $2n < f(x+1, n) - f(x, n)$.\\
La funzione minima è $f(i, n) = 2n \cdot i$.

Message complexity: 2 messaggi (sveglia e notifica) per link; ogni messaggio 1 bit (distinguiamo tra wakeup e notifica): $O(n)$ bits.\\
Dal vincolo: $t(\text{min}) + f(\text{min}, n) + n < 2n \cdot \text{min} + 2n \in O(n*\text{min})$ time units, dalla sostituzione di $f(i, n)$
con $2n \cdot i$ e dal fatto che $t(\text{min}) < n$

Attenzione: se il minimi sono ad esempio un indirizzo IP, il tempo di attesa minimo può diventare molto costoso.

\subsection{Universal Waitings}
Questo protocollo può essere esteso ad un grafo qualungue: ricezione del wakeup, attesa di $f(i, n)$ unità di tempo, notifica da parte del minimo.\\
Per generalizzarlo dobbiamo considerare il diametro del grafo come massima distanza tra due nodi al posto di $n-1$. \\

Come al solito, se le entità non hanno ID univoco, non riusciamo a fare election con protocolli deterministici\dots
quindi cerchiamo di fare l'analogo protocolli randomici.

Protocolli di:
\begin{itemize}
  \item \textbf{Monte Carlo}: terminano sempre ma il risultato è corretto con una probabilità che deve essere studiata.
  \item \textbf{Las Vegas}: il risultato è sempre corretto, se il protocollo termina (terminazione con una certa probabilità).
\end{itemize}

Las Vegas, nel caso della leader election, termina correttamente (se termina), ovvero con l'ID del leader.

Presupposto: non abbiamo gli identificativi \\
Idea: strategia del waiting ma con gli ID scelti randomicamente da ogni nodo. Facciamo dipendere il tempo di attesa dall'identificativo. \\
Idea: per aspettare il minor tempo possibile possiamo estrarre il numero casuale tra $0$ e $b$, con $b$ upper bound al massimo $N$. In realtà, può
essere anche più piccolo, l'importante è che l'ID più piccolo sia univoco, mentre sui valori più grandi possono ammassarsi anche più valori.

\begin{enumerate}
  \item Il protocollo è composto da round
  \item In ogni round ogni entità seleziona un intero casuale e lo pone uguale al suo identificativo
  \item Se il minimo è unico quella entità diventa il leader mediante il protocollo Waiting, altrimenti si procede con il prossimo round.
\end{enumerate}

Problema: come fa un nodo ad accorgersi se è leader o no? potrebbe succedere che due nodi estraggano lo stesso numero, guardare il numero
contenuto nel messaggio di notifica potrebbe non permettere di disambiguare. Siamo nel caso sincrono, possiamo usare il tempo che è
passato dall'invio della nostra notifica per determinare se è il nostro messaggio a tornare indietro o quello di un altro nodo.
Il nostro messaggio avrà atteso almeno $N$ tick prima di tornare.

Se così non è, inviamo un messaggio di \textbf{restart} a tutte le altre entità per iniziare un nuovo round.

Consideriamo un solo round.\\
Message Complexity: $O(n)$ bits per i messaggi di notifica e il restart.\\
Time Complexity: $O(n \cdot \text{min}_\text{id})$ dalla formula di waiting; dato che $\text{min}_\text{id} \leq b \rightarrow O(n \cdot b)$.
Se scegliamo $b=1$, il costo temporale diventa $O(n)$ per il singolo round.
Ogni nodo può estrarre 0 o 1. Se estrae 0 diventa il leader, in assenza di altri zero. \\
Quanti round? dobbiamo progettare l'estrazione del numero il modo che sia molto facie estrarre 1 e molto improbabile estrarre 0:
\begin{itemize}
  \item 0 con probabilità $\frac{1}{n}$
  \item 1 con probabilità $\frac{n-1}{n}$
\end{itemize}

Qual'è la prob che l'algoritmo termini dopo R round? \\
Scomponiamo il problema. Abbiamo vinto se un nodo sceglie zero e tutti gli altri 1; succede con prob $\frac{1}{n} (\frac{n-1}{n})^{n-1}$.
Qualunque degli $n$ nodi può scegliere lo zero\dots sommiamo le probabilità: $n \frac{1}{n} (\frac{n-1}{n})^{n-1} = (\frac{n-1}{n})^{n-1} \rightarrow_{n \rightarrow \inf} \frac{1}{e}$.

Per $n$ grandi (altrimenti non varrebbe l'approssimazione per $n$ a infinito). $P_\text{success} = \frac{1}{e} \quad P_\text{insuccess} = 1 - \frac{1}{e}$.
Prob of success at round R, meaning insucesses in previous $R-1$ rounds: $\frac{1}{e} (1 - \frac{1}{e})^{R-1}$.
Questa è la distribuzione geometrica. Valore atteso dopo $R$ round $e = 2.7 \approx 3$.

Tenendo conto della time complexity per un round. Final Time Complexity: $O(n)$.

Questo funziona se abbiamo un clock condiviso tra tutti i nodi. Significa avere i nodi su un sistema centralizzato, e.g. processi su una
stessa macchina. Tra nodi distribuiti servono politiche di allineamento degli orologi o serve l'utilizzo di un tempo logico.
